\cvheading{基于 LangGraph 的脑影像数据处理 Agent 与可恢复编排系统}{2026.05--至今}
{\small\color{OpenCurveMuted}Python / LangGraph / LLM API / Pydantic / Docker / SQLite / MySQL}\par
面向长耗时数据处理任务中命令配置复杂、失败恢复困难及结果缺少统一门禁的问题，设计并实现覆盖任务规划、受控执行、质量检查和结果持久化的 Agent 系统。
\cvitem{\textbf{受约束的任务规划：}由 LLM 将自然语言需求转换为结构化执行计划和受管参数，通过 Pydantic、Worker 注册表及步骤依赖规则限制可调用能力，禁止模型直接执行任意 Shell 或绕过 QC；在 30 条人工审核用例、每条 3 次的真实模型回放中，90/90 次结构化调用成功，原始计划与参数完全匹配率均为 100\%，禁用动作 0。}
\cvitem{\textbf{HITL 与质量门禁：}在计划确认、输入检查和输出 QC 阶段设置可恢复审批节点，将硬失败、待复核异常和人工决定写入审计状态；十名被试批处理中识别出 1 名掩膜/配准异常并阻断 MySQL 入库，后续根据输入检查调整预处理参数，使结构/功能掩膜 Dice 从 0.693 提升至 0.922。}
\cvitem{\textbf{恢复、复用与并行调度：}基于 SQLite checkpoint 与 thread ID 保存工作流状态，结合配置指纹、完成标记和文件锁实现幂等续跑及结果复用；真实批次复用 3 名被试并新计算 7 名，10/10 完成计算，另在 12 任务、5 次重复基准中实现 4 路并发相对串行 3.80 倍加速。}
\cvitem{\textbf{故障处理与持久化：}在 Worker 边界实现故障注入、暂时性错误重试和失败报告收敛，通过事务及幂等 upsert 保存获准任务；围绕计划约束、HITL、状态恢复、并发聚合、QC 阻断和数据库故障建立 63 项自动化测试，全部通过。}
